\documentclass[10pt,a4paper]{amsart}
\usepackage{amsmath,amssymb,amsthm,graphicx}
\newcommand{\RS}[1]{\textcolor{red}{#1}}

%\usepackage{refcheck}

\numberwithin{equation}{section}

\newtheorem{Theorem}{Theorem}[section]
\newtheorem*{Theorem*}{Theorem}
\newtheorem{Corollary}[Theorem]{Corollary}
\newtheorem{Lemma}[Theorem]{Lemma}
\newtheorem{Proposition}[Theorem]{Proposition}
 { \theoremstyle{definition}
\newtheorem{Definition}[Theorem]{Definition}
\newtheorem{Note}[Theorem]{Note}
\newtheorem{Example}[Theorem]{Example}
\newtheorem{Remark}[Theorem]{Remark} }

\newcommand{\sech}{\operatorname{sech}}
\usepackage{cleveref}
\crefname{originalequation}{original equation}{original equations}
\Crefname{originalequation}{Original equation}{Original equations}
\makeatletter
\newlabel{eqorig5_6}{{internal}{source}{}{sourcepointer-1}{}}
\newlabel{eqorig5_6@cref}{{[originalequation][0][]internal}{[1][0][]source}{}{}{}}
\makeatother
\crefname{originalequation}{original equation}{original equations}
\Crefname{originalequation}{Original equation}{Original equations}
\makeatletter
\newlabel{eq:btanalytic_v5_6}{{internal}{source}{}{sourcepointer-2}{}}
\newlabel{eq:btanalytic_v5_6@cref}{{[originalequation][0][]internal}{[1][0][]source}{}{}{}}
\makeatother
\crefname{originalequation}{original equation}{original equations}
\Crefname{originalequation}{Original equation}{Original equations}
\makeatletter
\newlabel{eqprime5_6}{{internal}{source}{}{sourcepointer-3}{}}
\newlabel{eqprime5_6@cref}{{[originalequation][0][]internal}{[1][0][]source}{}{}{}}
\makeatother
\title{Superposition formula for two Backlund transformations of the elliptic sinh-Gordon equation}
\author{Kelmer, Filipe, Tenenblat, Keti}
\date{doi:10.3842/SIGMA.2024.015; CC BY 4.0}
\begin{document}
\maketitle
\setcounter{section}{3}
\setcounter{Theorem}{1}
\noindent\textit{Editorial note.} All delivered mathematical text is the author's own, copied from the pinned SIGMA article source; no statement or proof step was rewritten, reordered or removed. The counted statement is the labeled theorem environment at source lines 642-660 and the counted proof is the proof environment at source lines 662-764.\par\medskip
\noindent\textit{Editorial note.} The retained preamble of this unit was generated mechanically for this delivery: it copies the author's own numbering, macro and theorem declaration block from the pinned source (source lines 9-27), with the publisher class and page setup lines removed.\par\medskip
\noindent\textit{Allowed standard prerequisites (editorial).}\begin{itemize}
\item The elliptic sinh-Gordon equation and its analytic Backlund transformations. same article, the corollary at source lines 330-346. Supplies the equation solved by the solution and the transformations relating it to its two companions.
\item The two solutions associated with the companions by the transformation parameters. same article, the displays at source lines 333, 339 and 343. Fix the equations satisfied by the companions and the general form of the transformation used in the elimination.
\end{itemize}\medskip
\ifdefined\hypertarget\hypertarget{sourcepointer-1}{}\fi
\noindent\textit{Source reference.} Original equation internal: sources/sigma-2024-015-superposition-formulae-for-the-geometric-b-2/pinned-original-source.tex, source lines 333--338. The counted statement or proof refers to that displayed equation satisfied by the companion solutions; it is the author's own earlier display or statement in the same article and the whole argument of the counted span is retained verbatim.\par\medskip
\ifdefined\hypertarget\hypertarget{sourcepointer-2}{}\fi
\noindent\textit{Source reference.} Original equation internal: sources/sigma-2024-015-superposition-formulae-for-the-geometric-b-2/pinned-original-source.tex, source lines 339--342. The counted statement or proof refers to that displayed analytic Backlund transformation used for this family; it is the author's own earlier display or statement in the same article and the whole argument of the counted span is retained verbatim.\par\medskip
\ifdefined\hypertarget\hypertarget{sourcepointer-3}{}\fi
\noindent\textit{Source reference.} Original equation internal: sources/sigma-2024-015-superposition-formulae-for-the-geometric-b-2/pinned-original-source.tex, source lines 343--346. The counted statement or proof refers to that displayed transformation equation for the companion of the second parameter; it is the author's own earlier display or statement in the same article and the whole argument of the counted span is retained verbatim.\par\medskip
\begin{Theorem}\label{th:super5}
Let $\alpha(x_1,x_2)$ be a solution of
\begin{equation}\label{elipsinh}
\alpha_{x_1x_1}+\alpha_{x_2x_2}=\sinh \alpha.
\end{equation}
Let $\alpha'$ and $\alpha''$ be solutions of the elliptic sine-Gordon equation associated to $\alpha$ by $\phi_1($resp.\ $\phi_2)$-B\"acklund transformation \eqref{eq:btanalytic_v5_6} where $r=0$ and $\tau=\pm 1$.
If $\phi_1\neq \phi_2$, then there exists a unique solution $\alpha^*$ of \eqref{elipsinh} which is associated to
$\alpha'$ and $\alpha''$ by $\phi_2($resp.\ $\phi_1)$-B\"acklund transformation~\eqref{eq:btanalytic_v5_6}, with $r=0$. Moreover, $\alpha^*$ is
given algebraically by
\begin{equation}\label{eq:superposcas5}
\tanh \frac{\alpha^*}{2}=
 \frac{P-Q \tanh \bigl({\frac{\alpha}{2}}\bigr)} {Q-P \tanh \bigl({\frac{\alpha}{2}}\bigr)},
\end{equation}
where
\begin{gather*}
P=\tau(\sinh\phi_2-\sinh\phi_1) \sin\biggl(\frac{\alpha''-\alpha'}{2}\biggr), \\
Q=(1+\sinh\phi_1\sinh\phi_2)\cos\biggl(\frac{\alpha''-\alpha'}{2}\biggr) -\cosh\phi_1\cosh\phi_2.
\end{gather*}
\end{Theorem}

\begin{proof} Let $\alpha$ be a solution of the elliptic sinh-Gordon equation, i.e., \eqref{eqorig5_6} with $r=0$.
Let $\alpha'$ and $\alpha''$ be solutions of the elliptic sine-Gordon equation, i.e., \eqref{eqprime5_6} with $r=0$, associated to $\alpha$
by $\phi_1$(resp.\ $\phi_2$)-B\"acklund transformation \eqref{eq:btanalytic_v5_6} with $r=0$, $\phi_1\neq \phi_2$. Then
 \begin{gather}
\alpha'_{x_1}-\alpha_{x_2}=\frac{2}{\cosh\phi_1}\sin \biggl(\frac{\alpha'}{2}\biggr)\cosh \Bigl(\frac{\alpha}{2}\Bigr)+\frac{2\tau\sinh\phi_1}{\cosh\phi_1} \cos\biggl(\frac{\alpha'}{2}\biggr)\sinh \Bigl(\frac{\alpha}{2}\Bigr),\nonumber\\
\alpha'_{x_2}+\alpha_{x_1}=-\frac{2}{\cosh\phi_1}\cos\biggl(\frac{\alpha'}{2}\biggr)\sinh\Bigl(\frac{\alpha}{2}\Bigr)+\frac{2\tau\sinh\phi_1}{\cosh\phi_1} \sin\biggl(\frac{\alpha'}{2}\biggr)\cosh\Bigl(\frac{\alpha}{2}\Bigr),\label{BTcase51}\\
\alpha''_{x_1}-\alpha_{x_2}=\frac{2}{\cosh\phi_2}\sin\biggl(\frac{\alpha''}{2}\biggr)\cosh \Bigl(\frac{\alpha}{2}\Bigr)+\frac{2\tau\sinh\phi_2}{\cosh\phi_2} \cos\biggl(\frac{\alpha''}{2}\biggr)\sinh\Bigl(\frac{\alpha}{2}\Bigr),\nonumber\\
\alpha''_{x_2}+\alpha_{x_1}=-\frac{2}{\cosh\phi_2}\cos\biggl(\frac{\alpha''}{2}\biggr)\sinh\Bigl(\frac{\alpha}{2}\Bigr)+\frac{2\tau\sinh\phi_2}{\cosh\phi_2} \sin\biggl(\frac{\alpha''}{2}\biggr)\cosh\Bigl(\frac{\alpha}{2}\Bigr).\label{BTcase52}
\end{gather}
Suppose $\alpha^*$ exists then it should satisfy the following
systems of equations
\begin{gather}
\alpha^*_{x_1}-\alpha'_{x_2}=\frac{2}{\cosh\phi_2}\sinh\biggl(\frac{\alpha^*}{2}\biggr)\cos\biggl(\frac{\alpha'}{2}\biggr)+\frac{2\tau\sinh\phi_2}{\cosh\phi_2} \cosh\biggl(\frac{\alpha^*}{2}\biggr)\sin\biggl(\frac{\alpha'}{2}\biggr),\nonumber\\
\alpha^*_{x_2}+\alpha'_{x_1}=-\frac{2}{\cosh\phi_ 2}\cosh\biggl(\frac{\alpha^*}{2}\biggr)\sin\biggl(\frac{\alpha'}{2}\biggr)+\frac{2\tau\sinh\phi_2}{\cosh\phi_2} \sinh\biggl(\frac{\alpha^*}{2}\biggr)\cos\biggl(\frac{\alpha'}{2}\biggr),\label{BTcase62}
\\
\alpha^*_{x_1}-\alpha''_{x_2}=\frac{2}{\cosh\phi_1}\sinh\biggl(\frac{\alpha^*}{2}\biggr)\cos\biggl(\frac{\alpha''}{2}\biggr)+\frac{2\tau\sinh\phi_1}{\cosh\phi_1} \cosh\biggl(\frac{\alpha^*}{2}\biggr)\sin\biggl(\frac{\alpha''}{2}\biggr),\nonumber\\
\alpha^*_{x_2}+\alpha''_{x_1}=-\frac{2}{\cosh\phi_1}\cosh\biggl(\frac{\alpha^*}{2}\biggr)\sin\biggl(\frac{\alpha''}{2}\biggr)+\frac{2\tau\sinh\phi_1}{\cosh\phi_1} \sinh\biggl(\frac{\alpha^*}{2}\biggr)\cos\biggl(\frac{\alpha''}{2}\biggr).\label{BTcase61}
\end{gather}
Subtracting the first (resp.\ second) equation of \eqref{BTcase52} from the first (resp.\ second) equation of~\eqref{BTcase51}, we get the system of equations
\begin{gather}
\alpha'_{x_1}-\alpha''_{x_1}=\frac{2}{\cosh\phi_1}\sin\biggl(\frac{\alpha'}{2}\biggr)\cosh \Bigl(\frac{\alpha}{2}\Bigr)+\frac{2\tau\sinh\phi_1}{\cosh\phi_1} \cos\biggl(\frac{\alpha'}{2}\biggr)\sinh\Bigl(\frac{\alpha}{2}\Bigr)
\nonumber\\ \hphantom{\alpha'_{x_1}-\alpha''_{x_1}=}{}
- \frac{2}{\cosh\phi_2}\sin\biggl(\frac{\alpha''}{2}\biggr)\cosh \Bigl(\frac{\alpha}{2}\Bigr)-\frac{2\tau\sinh\phi_2}{\cosh\phi_2} \cos\biggl(\frac{\alpha''}{2}\biggr)\sinh\Bigl(\frac{\alpha}{2}\Bigr),\nonumber\\
\alpha'_{x_2}-\alpha''_{x_2}= -\frac{2}{\cosh\phi_1}\cos\biggl(\frac{\alpha'}{2}\biggr)\sinh\Bigl(\frac{\alpha}{2}\Bigr)+\frac{2\tau\sinh\phi_1}{\cosh\phi_1} \sin\biggl(\frac{\alpha'}{2}\biggr)\cosh\Bigl(\frac{\alpha}{2}\Bigr)
\nonumber\\ \hphantom{\alpha'_{x_2}-\alpha''_{x_2}=}{}
+\frac{2}{\cosh\phi_2}\cos\biggl(\frac{\alpha''}{2}\biggr)\sinh\Bigl(\frac{\alpha}{2}\Bigr)-\frac{2\tau\sinh\phi_2}{\cosh\phi_2} \sin\biggl(\frac{\alpha''}{2}\biggr)\cosh\Bigl(\frac{\alpha}{2}\Bigr). \label{BT51minusBT52}
\end{gather}

Similarly, subtracting the first (resp.\ second) equation of \eqref{BTcase61} from the first (resp.\ second) equation of \eqref{BTcase62}, we get
\begin{gather}
\alpha'_{x_2}-\alpha''_{x_2}=-\frac{2}{\cosh\phi_2}\sinh\biggl(\frac{\alpha^*}{2}\biggr)\cos\biggl(\frac{\alpha'}{2}\biggr)-\frac{2\tau\sinh\phi_2}{\cosh\phi_2} \cosh\biggl(\frac{\alpha^*}{2}\biggr)\sin\biggl(\frac{\alpha'}{2}\biggr)
\nonumber\\ \hphantom{\alpha'_{x_2}-\alpha''_{x_2}=}{}
 +\frac{2}{\cosh\phi_1}\sinh\biggl(\frac{\alpha^*}{2}\biggr)\cos\biggl(\frac{\alpha''}{2}\biggr)+\frac{2\tau\sinh\phi_1}{\cosh\phi_1} \cosh\biggl(\frac{\alpha^*}{2}\biggr)\sin\biggl(\frac{\alpha''}{2}\biggr),\nonumber\\
\alpha'_{x_1}-\alpha''_{x_1}=-\frac{2}{\cosh\phi_ 2}\cosh\biggl(\frac{\alpha^*}{2}\biggr)\sin\biggl(\frac{\alpha'}{2}\biggr)+\frac{2\tau\sinh\phi_2}{\cosh\phi_2} \sinh\biggl(\frac{\alpha^*}{2}\biggr)\cos\biggl(\frac{\alpha'}{2}\biggr)
\nonumber\\ \hphantom{\alpha'_{x_1}-\alpha''_{x_1}=}{}
 +\frac{2}{\cosh\phi_1}\cosh\biggl(\frac{\alpha^*}{2}\biggr)\sin\biggl(\frac{\alpha''}{2}\biggr)- \frac{2\tau\sinh\phi_1}{\cosh\phi_1} \sinh\biggl(\frac{\alpha^*}{2}\biggr)\cos\biggl(\frac{\alpha''}{2}\biggr). \label{BT62minusBT61}
\end{gather}

Comparing equation \eqref{BT51minusBT52} with \eqref{BT62minusBT61}, we obtain two algebraic
expressions that must be satisfied by $\cosh(\alpha^*/2)$ and $\sinh(\alpha^*/2)$.
By solving this system of algebraic equations, we conclude that $\alpha^*/2$ is given by
\eqref{eq:superposcas5}.

Now we need to prove that $\alpha^*$ satisfies the elliptic sinh-Gordon equation \eqref{elipsinh}.
Recall that~$\alpha'$ and~$\alpha''$ are solutions of the elliptic sine-Gordon equation.
We introduce the following notation
\begin{gather}
A=\tau(\sinh\phi_2-\sinh\phi_1),\qquad B=\cosh\phi_1\cosh\phi_2,\qquad L=1+\sinh\phi_1\sinh\phi_2, \nonumber\\
 D=\alpha''-\alpha'.\label{ABLD}
\end{gather}
Then the functions $P$ and $Q$ are rewritten as
\begin{equation*}%\label{PQ}
P=A\sin\biggl(\frac{D}{2}\biggr),\qquad Q=L\cos\biggl(\frac{D}{2}\biggr)-B \qquad \mbox {and}\qquad B^2-L^2=A^2.
\end{equation*}
We denote the right-hand side of \eqref{eq:superposcas5} by
\begin{equation*}%\label{HG}
H=\frac{G-\tanh\bigl(\frac{\alpha}{2}\bigr)}{1-G\tanh\bigl(\frac{\alpha}{2}\bigr)}, \qquad \mbox{where}\quad G=\frac{P}{Q}.
\end{equation*}
It follows that proving that $\alpha^*$ satisfies \eqref{elipsinh} is equivalent to proving that
\begin{equation}\label{eqH}
\frac{2H}{1-H^2}|\nabla H|^2+ \triangle H-H=0.
\end{equation}
Since $\alpha$ satisfies \eqref{elipsinh} it follows that
\begin{gather*}
\triangle \tanh \Bigl(\frac{\alpha}{2}\Bigr)= \frac{1}{2\cosh^2\bigl(\frac{\alpha}{2}\bigr)}\Bigl(\sinh \alpha-\tanh\Bigl(\frac{\alpha}{2}\Bigr)|\nabla\alpha|^2\Bigr), \\
\Bigl|\nabla\tanh\Bigl(\frac{\alpha}{2}\Bigr)\Bigr|^2=\frac{1}{4\cosh^4 \bigl(\frac{\alpha}{2}\bigr)}
|\nabla \alpha|^2.
\end{gather*}
Therefore, proving \eqref{eqH} reduces to showing
that $G$ satisfies
\begin{equation}\label{eqG}
\frac{2G}{1-G^2}|\nabla G|^2+\triangle G-G(\cosh \alpha-G\sinh\alpha)=0.
\end{equation}

Long but straightforward computations show that
\begin{gather*}
 \triangle G = \frac{A}{2 Q^2}\biggl(L-B\cos\biggl(\frac{D}{2}\biggr) \biggr) {\triangle D}
-\frac{A}{4Q^3}\bigl(BQ+2A^2\bigr)\sin \biggl(\frac{D}{2}\biggr) |\nabla D|^2, \\
 |\nabla G|^2 = \frac{A^2}{4 Q^4}\biggl(L-B\cos\biggl(\frac{D}{2}\biggr) \biggr)^2 |\nabla D|^2,\\
 1-G^2 = \frac{1}{Q^2}\biggl(L-B\cos\biggl(\frac{D}{2}\biggr) \biggr)^2.
\end{gather*}
Moreover, since $D=\alpha''-\alpha'$ where $\alpha'$ and $\alpha''$ are solutions of the elliptic sine-Gordon equation, it follows that
\begin{eqnarray*}
 \triangle D=\sin \alpha''-\sin \alpha'= 2\sin\frac{D}{2}\cos\biggl(\frac{\alpha''+\alpha'}{2}\biggr).
\end{eqnarray*}
Hence proving \eqref{eqG} reduces to proving
\begin{equation}\label{eqD}
-B|\nabla D|^2+\frac{2}{\sin\bigl(\frac{D}{2}\bigr)}\biggl(L-B\cos\biggl(\frac{D}{2}\biggr)\biggr)\triangle D-4Q\cosh
\alpha+4A\sin\biggl(\frac{D}{2}\biggr)\sinh\alpha=0.
\end{equation}

Now $\alpha'$ and $\alpha''$ satisfy \eqref{BTcase51} and \eqref{BTcase52} respectively.
These systems of equations provide
\begin{gather*}
 |\nabla D|^2= \frac{4}{B}
 \biggl[\cosh \alpha \biggl( B-L\cos\biggl(\frac{D}{2}\biggr)\biggr) +
 \cos \biggl(\frac{\alpha''+\alpha'}{2}\biggr)\biggl( L-B\cos\biggl(\frac{D}{2}\biggr)\biggr)
 \nonumber \\ \hphantom{|\nabla D|^2= \frac{4}{B}\biggl[}{}
 + A\sin \biggl(\frac{D}{2}\biggr)\sinh \alpha\biggr].
\end{gather*}
This last expression proves that \eqref{eqD} is identically satisfied, which concludes the proof that
$\alpha^*$ satisfies the elliptic sinh-Gordon equation
\end{proof}

\end{document}
